\section{Búsqueda en el momento de la inferencia: de la ejecución de una sola trayectoria a la decisión con posibilidad de retroceso}

\subsection{Por qué es necesario introducir la búsqueda}
La conclusión que da el expositor es muy clara: sin capacidad de búsqueda, en cuanto el agente se desvía es muy difícil que regrese a la trayectoria correcta. Esto es especialmente cierto en escenarios web, donde una acción errónea cambia el estado de la página, el espacio de operaciones posteriores cambia en consecuencia, y el costo de recuperación es muy alto.

\begin{importantbox}{Beneficios directos de la búsqueda}
\begin{itemize}
    \item Permite evaluar en paralelo varias trayectorias candidatas;
    \item Permite el retroceso y reintentar por ramas;
    \item Convierte «acertar de un solo intento» en «elegir la mejor tras varios intentos».
\end{itemize}
\end{importantbox}

\subsection{Estrategia básica: Rejection Sampling}
El método más simple consiste en generar primero varias trayectorias y luego filtrarlas con un criterio de resultado final. Es sencillo de implementar y favorece la paralelización, pero no es muy eficiente, porque una gran cantidad de trayectorias inválidas desperdicia el presupuesto de inferencia.

\begin{knowledgebox}{Cuándo funciona Rejection Sampling}
Cuando la tarea es corta, el verificador es confiable y hay suficientes recursos para paralelizar, este método puede servir como referencia de bajo costo de ingeniería y aportar mejoras rápidas.
\end{knowledgebox}

\subsection{Estrategia avanzada: guía por valor y búsqueda en árbol}
Cuando aumenta la longitud de la tarea, el muestreo puro deja de funcionar rápidamente. En ese momento hay que introducir una función de valor o un Judge que estime el potencial de los estados intermedios, dando prioridad a expandir los nodos de mayor valor, para lograr «explorar mientras se podan ramas».

\begin{figure}[H]
\centering
\includegraphics[width=0.82\textwidth]{frame-07-search-eval.jpg}
\caption{Video en aproximadamente 00:34:39: el expositor explica cómo se evalúa el resultado de una trayectoria en una rama de búsqueda y se decide si continuar o terminar}
\end{figure}

\begin{table}[H]
\centering
\small
\begin{tabular}{p{2.8cm}p{3.8cm}p{4.4cm}}
\toprule
\textbf{Método} & \textbf{Ventajas} & \textbf{Desventajas} \\
\midrule
Ejecución de una sola trayectoria & Baja latencia, implementación sencilla & Los errores no son recuperables, la tasa de éxito depende mucho del ruido de una sola decisión \\
Rejection Sampling & Favorece la paralelización, fácil de poner en producción & Desperdicio de cómputo notable, la ganancia disminuye en tareas largas \\
Búsqueda guiada por valor & Alta eficiencia de muestreo, permite retroceso & Ingeniería compleja, sensible a la calidad del Judge \\
\bottomrule
\end{tabular}
\caption{Comparación de estrategias de razonamiento del agente}
\end{table}

\begin{warningbox}{La búsqueda no es gratuita}
La búsqueda aumenta de manera notable el costo de inferencia. Si no hay un buen mecanismo de poda y reutilización de nodos, el sistema puede caer en un estado de «presupuesto agotado sin haber convergido».
\end{warningbox}

\subsection{Resumen de la sección}
En las tareas web de cadena larga, la capacidad de búsqueda no es opcional, sino la fuente central de la robustez. Lo importante no es «si se busca», sino «cómo se busca de manera eficiente».
